% TEMPLATE-MILC-v3.tex - plantilla maestra para todas las guias MILC v3
% Ver scripts/generadores/build-guia-milc.py para la lista completa de
% claves esperadas. El script valida que no queden placeholders sin
% reemplazar antes de escribir el archivo final.

\documentclass[11pt,letterpaper]{article}
\usepackage[margin=1.75cm]{geometry}
\usepackage[table]{xcolor}
\usepackage{fontspec}
\usepackage[spanish]{babel}
\usepackage{tikz}
\usepackage[most]{tcolorbox}
\usepackage{tabularx}
\usepackage{array}
\usepackage{fancyhdr}
\usepackage{titlesec}
\usepackage{ragged2e}
\usepackage{enumitem}
\usepackage{microtype}

\setmainfont{Helvetica Neue}
\setsansfont{Helvetica Neue}
\setlength{\parindent}{0pt}
\setlength{\parskip}{6pt}
\hyphenpenalty=200
\exhyphenpenalty=200
\pretolerance=400
\tolerance=2000
\setlength{\emergencystretch}{1em}
\renewcommand{\arraystretch}{1.25}
\setlist{leftmargin=5mm,itemsep=1mm,topsep=1mm}
\newcolumntype{Y}{>{\RaggedRight\arraybackslash}X}

\definecolor{milcMagenta}{HTML}{E6007E}
\definecolor{milcVino}{HTML}{5A0038}
\definecolor{milcTurquesa}{HTML}{0093A5}
\definecolor{milcVerde}{HTML}{58A923}
\definecolor{milcMostaza}{HTML}{E5B400}
\definecolor{milcOcre}{HTML}{B86600}
\definecolor{milcNegro}{HTML}{241816}
\definecolor{milcGris}{HTML}{F7F7F4}
\definecolor{milcLinea}{HTML}{E8E2DD}
\definecolor{milcGradePrimary}{HTML}{84CC16}
\definecolor{milcGradeDark}{HTML}{4D7C0F}
\definecolor{milcGradeSoft}{HTML}{ECFCCB}
\definecolor{milcGradeAccent}{HTML}{CFE7FF}
\definecolor{milcGradeText}{HTML}{FFFFFF}
\color{milcNegro}

\pagestyle{fancy}
\fancyhf{}
\lhead{\small Tecnología e Informática · Grado 7}
\rhead{\small Guía 26 · MILC v3}
\cfoot{\small\thepage}
\renewcommand{\headrulewidth}{0pt}

\newtcolorbox{titlebox}[1]{
  enhanced,colback=#1,colframe=#1,arc=7mm,boxrule=0pt,
  left=7mm,right=7mm,top=3mm,bottom=3mm,
  before skip=8pt,after skip=8pt,
  fontupper=\sffamily\bfseries\Large\color{white}
}

\newtcolorbox{softbox}[3]{
  enhanced,colback=#2,colframe=#1,borderline west={4pt}{0pt}{#1},
  arc=2mm,boxrule=.4pt,left=4mm,right=4mm,top=3mm,bottom=3mm,
  before skip=6pt,after skip=8pt,title={#3},coltitle=white,
  colbacktitle=#1,fonttitle=\sffamily\bfseries,fontupper=\RaggedRight,
  attach boxed title to top left={xshift=3mm,yshift=-2mm},
  boxed title style={arc=2mm, boxrule=0pt}
}

\newtcolorbox{drawbox}[2]{
  enhanced,colback=white,colframe=#1,arc=4mm,boxrule=1pt,
  left=5mm,right=5mm,top=4mm,bottom=4mm,before skip=8pt,after skip=8pt,
  title={#2},coltitle=white,colbacktitle=#1,fonttitle=\sffamily\bfseries,
  fontupper=\RaggedRight,
  attach boxed title to top left={xshift=4mm,yshift=-2mm},
  boxed title style={arc=3mm, boxrule=0pt}
}

\newtcolorbox{infoband}[1]{
  enhanced,colback=#1!8,colframe=#1!50,arc=2mm,boxrule=.5pt,
  left=4mm,right=4mm,top=2mm,bottom=2mm,before skip=4pt,after skip=4pt,
  fontupper=\small\RaggedRight
}

% Puente narrativo entre secciones (contrato editorial Conéctate).
% Da continuidad: "Acabas de X. Ahora vas a Y porque Z."
% Visualmente: banda izquierda en color del grado, texto en itálica pequeña.
\newtcolorbox{puentebox}{
  enhanced,colback=milcGradeSoft,colframe=milcGradeDark,
  borderline west={3pt}{0pt}{milcGradeDark},
  arc=0mm,boxrule=0pt,left=5mm,right=4mm,top=2.5mm,bottom=2.5mm,
  before skip=4pt,after skip=8pt,
  fontupper=\itshape\small\color{milcNegro}\RaggedRight
}
\newcommand{\puente}[1]{%
  \begin{puentebox}%
    \textcolor{milcGradeDark}{\textbf{\textrightarrow}}\hspace{2mm}#1%
  \end{puentebox}%
}

\newcommand{\linead}[1]{\noindent\color{milcLinea}\rule{#1}{0.5pt}\color{milcNegro}}
\newcommand{\respuesta}[1]{\par\vspace{2mm}\foreach \n in {1,...,#1}{\noindent\color{milcLinea}\rule{\linewidth}{0.5pt}\par\vspace{5mm}}\color{milcNegro}}
\newcommand{\checkbox}{\textcolor{milcGradeDark}{\fbox{\phantom{x}}}\hspace{5pt}}

\newcommand{\phasecard}[4]{
\begin{tcolorbox}[enhanced,colback=#3,colframe=#2,arc=3mm,boxrule=.5pt,height=3.05cm,valign=center,halign=center,left=1.5mm,right=1.5mm,top=2mm,bottom=2mm]
{\sffamily\bfseries\fontsize{22}{24}\selectfont\textcolor{#2}{#1}}\par
{\sffamily\bfseries\small #4}
\end{tcolorbox}
}

\begin{document}
\thispagestyle{empty}
\begin{tikzpicture}[remember picture,overlay]
  \fill[white] (current page.south west) rectangle (current page.north east);
  \path[fill=milcGradePrimary,rounded corners=16mm]
    ([xshift=.55cm,yshift=-.55cm]current page.north west) rectangle
    ([xshift=-.55cm,yshift=3.65cm]current page.south east);

  \node[anchor=north west,text=milcGradeText] at ([xshift=2.0cm,yshift=-1.85cm]current page.north west)
    {\sffamily\bfseries\fontsize{14}{16}\selectfont GRADO};
  \node[anchor=north west,text=milcGradeText,opacity=.82] at ([xshift=2.0cm,yshift=-2.75cm]current page.north west)
    {\sffamily\bfseries\fontsize{9}{11}\selectfont SERIE GUÍAS · TECNOLOGÍA E INFORMÁTICA};

  \begin{scope}[shift={([xshift=-3.05cm,yshift=-2.55cm]current page.north east)}]
    \fill[white,opacity=.80] (-.62,-.52) rectangle (.62,.52);
    \draw[milcGradeText,line width=1pt] (-.62,-.52) rectangle (.62,.52);
    \fill[milcGradeDark] (-.42,-.38) rectangle (-.22,.06);
    \fill[milcGradeAccent] (-.10,-.38) rectangle (.10,.24);
    \fill[milcGradePrimary!60!black] (.22,-.38) rectangle (.42,.38);
  \end{scope}

  \node[anchor=west,text=milcGradeText] at ([xshift=1.9cm,yshift=-12.85cm]current page.north west)
    {\sffamily\bfseries\fontsize{124}{124}\selectfont 7};
  \draw[milcGradeText,line width=1.7pt]
    ([xshift=4.52cm,yshift=-11.68cm]current page.north west) circle (.13cm);

  \node[anchor=west,text=milcGradeText,text width=8.7cm,align=left] at ([xshift=10.35cm,yshift=-11.6cm]current page.north west)
    {
     {\sffamily\bfseries\fontsize{19}{22}\selectfont Prompting\\básico ---\\cómo hablarle\\a la IA\\para que\\responda bien}\\[4mm]
     {\sffamily\bfseries\fontsize{23}{26}\selectfont Guía 26-7-TIC}\\[2mm]
     {\sffamily\fontsize{13}{16}\selectfont Período 3 · Inteligencia Artificial}\\[5mm]
     {\sffamily\bfseries\fontsize{11}{13}\selectfont Institución Educativa Sor María Juliana}\\[1mm]
     {\sffamily\fontsize{10}{12}\selectfont Autor: Dr. Álvaro Cárdenas Orozco}
    };

  \path[fill=milcGradeSoft,rounded corners=7mm]
    ([xshift=1.35cm,yshift=.85cm]current page.south west) rectangle
    ([xshift=-1.35cm,yshift=4.25cm]current page.south east);
  \draw[milcGradeDark,line width=.65pt,rounded corners=7mm]
    ([xshift=1.35cm,yshift=.85cm]current page.south west) rectangle
    ([xshift=-1.35cm,yshift=4.25cm]current page.south east);
  \node[anchor=center,text=milcNegro,text width=17.0cm,align=center] at ([yshift=2.55cm]current page.south)
    {\sffamily\fontsize{10}{12}\selectfont Metodología MILC · Apertura ancestral · Pensamiento computacional · Triángulo Dussel-Estoicismo-Floridi\\
    \textcolor{milcGradeDark}{\bfseries Producto:} En el cuaderno: \textbf{(a) tabla de 10 prompts} (5 columnas: \emph{prompt malo} / \emph{problemas que tiene} / \emph{prompt mejorado con fórmula CTRF} / \emph{respuesta obtenida del LLM} / \emph{¿qué mejoró?}); \textbf{(b) tu fórmula CTRF anotada} (Contexto + Tarea + Restricciones + Formato) con ejemplo personal; \textbf{(c) un prompt complejo} para una tarea real que tengas que hacer esta semana. Firma + reflexión.};
\end{tikzpicture}
\mbox{}
\newpage

\section*{Apertura: Prompting básico --- cómo hablarle a la IA para que te responda bien}

\begin{softbox}{milcGradeDark}{milcGradeSoft}{Saber ancestral}
\textbf{Doña Mercedes la maestra rural del Valle del Cauca enseñaba a los niños a preguntarle bien al consejero del barrio.} \emph{``No vayan donde don Lucho con un `dígame algo' al aire''}, les decía. \emph{``Vayan con la pregunta concreta: `Don Lucho, mi papá quiere vender la finca, ¿qué debería considerar antes de decidir?' Así él les responde algo útil''}. Doña Mercedes sabía que \textbf{la calidad de la respuesta depende de la calidad de la pregunta}. Si llegabas vago, recibías consejo vago. Si llegabas específico, recibías consejo específico. Aprender a \emph{formular bien una pregunta} era habilidad que se enseñaba desde niño: \emph{``contextualicen su situación, sean claros, díganle qué necesitan exactamente''}. \textbf{Hoy con la IA pasa lo mismo}. Un \emph{prompt} (pregunta a la IA) bien formulado te da respuesta valiosa. Un prompt vago te da basura. \emph{Aprender a preguntar bien es 50\% del oficio de usar IA con criterio}.
\end{softbox}

\begin{softbox}{milcTurquesa}{milcGris}{Saber contemporáneo}
Un \textbf{prompt} es la instrucción o pregunta que le das a una IA. La habilidad de formular buenos prompts se llama \textbf{prompting} y es la habilidad más rentable del momento (incluso hay personas que trabajan como \emph{prompt engineers}, ganando muy bien). \textbf{La fórmula básica CTRF tiene 4 partes:} \textbf{(1) Contexto:} quién eres y por qué preguntas. \emph{``Soy estudiante de 7° grado en Colombia''}. \textbf{(2) Tarea:} qué quieres que haga la IA. \emph{``Explícame el teorema de Pitágoras''}. \textbf{(3) Restricciones:} cómo debe ser la respuesta. \emph{``Usa lenguaje sencillo, con un ejemplo del fútbol''}. \textbf{(4) Formato:} en qué forma. \emph{``Hazlo en 3 párrafos cortos''}. Cuando aplicas las 4 partes, transformas un prompt vago como \emph{``explícame Pitágoras''} en uno que da respuesta personalizada y útil. \textbf{Esa diferencia es la diferencia entre IA útil y IA basura}.
\end{softbox}

\begin{softbox}{milcMostaza}{milcGris}{Pregunta-puente}
\textit{Si tu profe te dice \emph{``dime algo''} sin contexto ni especifico, ¿\textbf{qué le respondes}? ¿Y si te dice \emph{``cuéntame 3 cosas que aprendiste hoy en clase de mate''}? La diferencia es prompting --- así pasa también con la IA.}
\end{softbox}

\begin{softbox}{milcVerde}{milcGris}{Saber y saber hacer}
Al terminar la clase: (1) podrás \textbf{identificar} prompts buenos vs malos; (2) sabrás \textbf{aplicar} la fórmula CTRF (Contexto + Tarea + Restricciones + Formato); (3) podrás \textbf{evaluar} la mejora entre un prompt antes y uno después; (4) habrás \textbf{creado} 10 prompts mejorados probados con un LLM real.
\end{softbox}

\small
\begin{tabularx}{\textwidth}{>{\bfseries}p{3.0cm}Y}
\rowcolor{milcGradePrimary!12}Autor & Dr. Álvaro Cárdenas Orozco\\
DBA & Reconoce los fundamentos, usos y límites éticos de la Inteligencia Artificial (MEN, T\&I 7°)\\
Referentes & Saber ancestral del consejero · Modelos de lenguaje · Ética algorítmica y prompting\\
Producto final & En el cuaderno: \textbf{(a) tabla de 10 prompts} (5 columnas: \emph{prompt malo} / \emph{problemas que tiene} / \emph{prompt mejorado con fórmula CTRF} / \emph{respuesta obtenida del LLM} / \emph{¿qué mejoró?}); \textbf{(b) tu fórmula CTRF anotada} (Contexto + Tarea + Restricciones + Formato) con ejemplo personal; \textbf{(c) un prompt complejo} para una tarea real que tengas que hacer esta semana. Firma + reflexión.\\
\end{tabularx}
\normalsize

\puente{Hoy haces 4 cosas: distingues prompts buenos vs malos, aprendes la fórmula CTRF, mejoras 10 prompts, los pruebas con LLM.}

\begin{titlebox}{milcGradeDark}
Ruta MILC: así vamos a aprender
\end{titlebox}
\begin{tabularx}{\textwidth}{>{\centering\arraybackslash}X>{\centering\arraybackslash}X>{\centering\arraybackslash}X>{\centering\arraybackslash}X}
\phasecard{1}{milcVerde}{milcGris}{Escuta\\[-1mm]\small Observo, escucho y pregunto.} &
\phasecard{2}{milcTurquesa}{milcGris}{Sistematización\\[-1mm]\small Pensamiento computacional.} &
\phasecard{3}{milcMagenta}{milcGris}{Praxis\\[-1mm]\small Construyo y aplico.} &
\phasecard{4}{milcVino}{milcGris}{Evaluación\\[-1mm]\small Reflexión liberadora.}
\end{tabularx}


\puente{Empezamos analizando 6 prompts (3 buenos, 3 malos) para reconocer la diferencia.}

\begin{titlebox}{milcVerde}
Fase 1 · Escuta
\end{titlebox}

\begin{softbox}{milcVerde}{milcGris}{Escena para escuchar}
\textbf{Actividad 1 · ANALIZA --- 6 prompts: ¿bueno o malo?} (10 min · individual). Lee estos 6 prompts y marca en el cuaderno BUENO (B) o MALO (M), con 1 línea explicando por qué. \emph{(1) ``Dime algo''}. \emph{(2) ``Soy estudiante de 7° grado. Necesito que me expliques el ciclo del agua para una tarea de Ciencias. Quiero la explicación en lenguaje sencillo, con un ejemplo del río Cauca, en 3 párrafos''}. \emph{(3) ``Bolívar''}. \emph{(4) ``Soy estudiante en Colombia. Hazme un resumen de los 5 logros más importantes de Simón Bolívar para América Latina, en máximo 200 palabras''}. \emph{(5) ``Ayúdame''}. \emph{(6) ``Soy estudiante de 7° grado. Quiero un correo formal al rector pidiéndole permiso para ir al baño durante un examen. 4 líneas máximo, con saludo y despedida formal''}.
\end{softbox}

\checkbox Leí los 6 prompts.\\
\checkbox Marqué B o M en cada uno.\\
\checkbox Identifiqué por qué los malos son malos.

\begin{infoband}{milcVerde}
\textbf{Lo que va al cuaderno (Actividad 1 · ANALIZA).} Título: «Actividad 1 --- 6 prompts: bueno o malo». \textbf{Formato:} lista numerada con B/M + 1 línea de explicación. \textbf{Extensión:} 6 líneas. \textbf{Sabes que terminaste cuando} reconoces que los buenos tienen contexto + tarea + restricciones + formato.
\end{infoband}

\puente{Después aprendes la fórmula CTRF en detalle.}

\begin{titlebox}{milcTurquesa}
Fase 2 · Sistematización con pensamiento computacional
\end{titlebox}

\begin{softbox}{milcTurquesa}{milcGris}{Cuatro pilares aplicados al tema}
Los 6 prompts se distinguen por una cosa: los buenos tienen \textbf{la fórmula CTRF aplicada}. Ahora aprende la fórmula en detalle + cómo aplicarla.
\end{softbox}

\small
\begin{tabularx}{\textwidth}{>{\bfseries\RaggedRight\arraybackslash}p{3.6cm}Y}
\rowcolor{milcTurquesa!90}\color{white}Pilar & \color{white}\bfseries Aplicación al tema\\
1. Descomposición & \textbf{Parte 1 de CTRF --- CONTEXTO.} Dile a la IA quién eres y por qué preguntas. \emph{Ejemplos buenos:} \emph{``Soy estudiante de 7° grado en Colombia''}; \emph{``Soy profesora de matemáticas en secundaria''}; \emph{``Tengo 12 años y vivo en Cartago, Valle del Cauca''}. \textbf{¿Por qué importa?} La IA adapta el tono y la profundidad según tu contexto. Sin contexto, te responde como si fueras adulto experto (puede ser muy técnico). \textbf{Parte 2 --- TAREA.} Qué quieres que haga, con verbo claro al inicio. \emph{Ejemplos:} \emph{``Explícame...''}, \emph{``Resúmeme...''}, \emph{``Tradúceme...''}, \emph{``Compárame...''}, \emph{``Lista 5 ideas para...''}. \textbf{Sin tarea, la IA no sabe qué quieres}.\\
2. Reconocimiento de patrones & \textbf{Parte 3 --- RESTRICCIONES.} Cómo debe ser la respuesta. \emph{Restricciones comunes:} \emph{lenguaje sencillo / nivel para 7° grado / con un ejemplo del Valle del Cauca / sin jerga técnica / contextualizado a Colombia / con palabras simples}. \textbf{Sin restricciones, la respuesta puede ser muy genérica o demasiado avanzada}. \textbf{Parte 4 --- FORMATO.} En qué forma quieres la respuesta. \emph{Formatos comunes:} \emph{en 3 párrafos / en lista numerada / en tabla / en máximo 200 palabras / como un correo formal con saludo y despedida / como un resumen con bullets}. \textbf{El formato hace la respuesta usable directamente}.\\
3. Abstracción & \textbf{Aplicación de CTRF: un prompt completo.} \emph{Prompt malo:} \emph{``Pitágoras''}. \emph{Prompt CTRF aplicado:} \emph{``[CONTEXTO] Soy estudiante de 7° grado en Cartago, Valle del Cauca. [TAREA] Explícame el teorema de Pitágoras. [RESTRICCIONES] Usa lenguaje sencillo, con un ejemplo del fútbol o de algo cotidiano. [FORMATO] En 3 párrafos cortos, máximo 150 palabras''}. Compara las 2 versiones: el segundo le dice al LLM EXACTAMENTE qué hacer.\\
4. Algoritmo conceptual & \textbf{La técnica de iteración.} A veces el primer prompt no da respuesta perfecta. La técnica adulta: \textbf{itera}. \emph{Si la respuesta es muy larga:} \emph{``Resúmela en 3 líneas''}. \emph{Si es muy técnica:} \emph{``Explícame eso como si fuera para un niño de 10''}. \emph{Si falta ejemplo:} \emph{``Dame un ejemplo concreto del Valle del Cauca''}. La conversación con LLM es ida y vuelta, no monólogo. Quien itera obtiene respuestas mejores que quien acepta la primera.\\
\end{tabularx}
\normalsize

\begin{drawbox}{milcTurquesa}{Fórmula CTRF + iteración}
\textbf{CTRF:} \\[1mm]
\textbf{C} - CONTEXTO: quién eres, por qué preguntas. \\[1mm]
\textbf{T} - TAREA: qué quieres que haga (verbo claro). \\[1mm]
\textbf{R} - RESTRICCIONES: cómo debe ser la respuesta. \\[1mm]
\textbf{F} - FORMATO: en qué forma. \\[1mm]
\textbf{Iteración:} si no es perfecta, pide ajustes (más corto, más simple, con ejemplo).
\end{drawbox}

\begin{softbox}{milcMostaza}{milcGris}{Errores comunes y su corrección}
\textbf{Errores frecuentes al hacer prompts.} (1) \emph{Prompt de 1-2 palabras:} \emph{``Bolívar''}. La IA no sabe qué quieres. (2) \emph{Sin contexto:} la IA asume que eres experto y te responde técnico. (3) \emph{Sin restricciones:} respuesta genérica. (4) \emph{Sin formato:} respuesta en bloque difícil de usar. (5) \emph{Aceptar primera respuesta sin iterar:} pierdes la oportunidad de afinar. (6) \emph{Querer respuesta perfecta de primer intento:} usualmente toma 2-3 iteraciones.
\end{softbox}

\begin{infoband}{milcTurquesa}
\textbf{Lo que va al cuaderno (Actividad 2 · EXPLICA).} Título: «Actividad 2 --- Fórmula CTRF». \textbf{Formato:} 4 bloques (C, T, R, F) con 1 ejemplo de cada uno + ejemplo de prompt completo CTRF aplicado. \textbf{Extensión:} 10 líneas. \textbf{Sabes que terminaste cuando} podrías escribir un prompt CTRF en 30 segundos para cualquier necesidad.
\end{infoband}

\puente{Con todo claro, mejoras 10 prompts tuyos y los pruebas con un LLM.}

\begin{titlebox}{milcMagenta}
Fase 3 · Praxis con pensamiento computacional
\end{titlebox}

\begin{softbox}{milcMagenta}{milcGris}{Construyendo el producto con los cuatro pilares}
\textbf{Actividad 3 · APLICA --- 10 prompts mejorados probados con LLM} (32 min · individual). Vas a tomar 10 prompts vagos y mejorarlos con CTRF. Después pruebas al menos 5 en un LLM real y comparas resultados.
\end{softbox}

\small
\begin{tabularx}{\textwidth}{>{\bfseries\RaggedRight\arraybackslash}p{3.6cm}Y}
\rowcolor{milcMagenta!90}\color{white}Pilar & \color{white}\bfseries Aplicación al producto\\
Descomposición & \textbf{Paso 1 --- Tabla de 10 prompts.} En una hoja del cuaderno, dibuja tabla de \textbf{10 filas y 5 columnas}. Columnas: \emph{(1) Prompt original (malo)}, \emph{(2) Problemas que tiene}, \emph{(3) Prompt mejorado con CTRF}, \emph{(4) Respuesta del LLM (resumen)}, \emph{(5) ¿Qué mejoró?}.\\
Patrones & \textbf{Paso 2 --- Llena columna 1 con 10 prompts originales (malos).} Inventa o copia de tu cabeza. Sugerencias: \emph{(1) ``Bolívar''}, \emph{(2) ``Pitágoras''}, \emph{(3) ``Cuéntame de Cartago''}, \emph{(4) ``Recomiéndame libros''}, \emph{(5) ``Corrígeme este texto''}, \emph{(6) ``Hazme una tarea''}, \emph{(7) ``Tradúcelo''}, \emph{(8) ``Explícame internet''}, \emph{(9) ``Hazme un poema''}, \emph{(10) ``Tabla periódica''}.\\
Abstracción & \textbf{Paso 3 --- Columna 2: problemas.} Para cada prompt malo, anota en 1 línea qué falta: \emph{``Falta contexto y tarea específica''}, \emph{``No dice formato ni nivel''}, etc. Esto te entrena la mente en detectar los huecos.\\
Algoritmo de armado & \textbf{Paso 4 --- Columna 3: prompts mejorados con CTRF.} Reescribe cada uno aplicando la fórmula. Ejemplo para el (1): \emph{``Soy estudiante de 7° grado en Colombia (C). Resúmeme los 3 logros principales de Simón Bolívar (T). Usa lenguaje sencillo, contextualizado a un colombiano (R). Hazlo en máximo 100 palabras (F)''}.\\
\end{tabularx}
\normalsize

\begin{drawbox}{milcMagenta}{Antes de cerrar la S6}
\checkbox La tabla tiene 10 prompts con sus 5 columnas. \\[1mm]
\checkbox Cada prompt mejorado tiene las 4 partes de CTRF. \\[1mm]
\checkbox Probé mínimo 5 prompts con un LLM real. \\[1mm]
\checkbox Comparé respuesta antes/después en columna 5. \\[1mm]
\checkbox Tengo 1 prompt complejo para una tarea real. \\[1mm]
\checkbox La reflexión está firmada.
\end{drawbox}

\begin{softbox}{milcMagenta}{milcGris}{Plantilla de trabajo}
\textbf{Estructura.} \textit{Paso 1:} tabla 10 × 5. \textit{Paso 2:} 10 prompts originales. \textit{Paso 3:} problemas. \textit{Paso 4:} CTRF aplicado. \textit{Paso 5:} pruebas con LLM. \textit{Paso 6:} reflexión + prompt para tarea real + firma.
\end{softbox}

\begin{infoband}{milcMagenta}
\textbf{Lo que va al cuaderno (Actividad 3 · APLICA).} Título: «Actividad 3 --- 10 prompts mejorados con CTRF». \textbf{Formato:} tabla 10x5 + reflexión + prompt para tarea real. \textbf{Extensión:} 1 página y media. \textbf{Sabes que terminaste cuando} puedes escribir cualquier prompt en formato CTRF.
\end{infoband}

\puente{El producto es la tabla de 10 prompts + fórmula CTRF + reflexión.}

\begin{titlebox}{milcMostaza}
Producto de la sesión
\end{titlebox}
\textbf{10 prompts mejorados con CTRF · probados con LLM · firmado}.
\begin{softbox}{milcMostaza}{milcGris}{Criterios de calidad}
(1) La tabla tiene 10 prompts con sus 5 columnas completas. (2) Cada prompt mejorado tiene las 4 partes de CTRF identificables. (3) Probé mínimo 5 con un LLM real. (4) Hay análisis comparativo de antes/después. (5) Hay 1 prompt para tarea real esta semana.
\end{softbox}

\puente{Antes de salir, verifica que tus prompts mejorados contienen las 4 partes.}

\begin{titlebox}{milcVino}
Fase 4 · Evaluación liberadora — cinco dimensiones
\end{titlebox}

\begin{softbox}{milcVino}{milcGris}{Sentido de la evaluación}
Evaluar no es solo revisar una nota. Es reconocer qué aprendí, qué cambió en mí, qué emociones manejé, cómo me relaciono con la ciudad y la comunidad, y qué le dejaré a quien viene detrás.
\end{softbox}

\textbf{Mi reflexión liberadora en cinco dimensiones.} Completa cada frase iniciada:

\small
\begin{tabularx}{\textwidth}{>{\bfseries\RaggedRight\arraybackslash}p{3.4cm}Y>{\RaggedRight\arraybackslash}p{4.6cm}}
\rowcolor{milcVino}\color{white}Dimensión & \color{white}\bfseries Mi respuesta (frame starter) & \color{white}\bfseries Algo que descubrí\\
1. Personal & Lo que más cambió en mí fue \linead{6.5cm} & \linead{4.2cm}\\
2. Emocional & La emoción más fuerte fue \linead{2.5cm} y la regulé \linead{3cm} & \linead{4.2cm}\\
3. Ciudadana & En democracia esto importa porque \linead{6cm} & \linead{4.2cm}\\
4. Local (Cartago) & En mi barrio sería útil para \linead{6.5cm} & \linead{4.2cm}\\
5. Intergeneracional & A mi abuela/abuelo le diría \linead{6.5cm} & \linead{4.2cm}\\
\end{tabularx}
\normalsize

\begin{infoband}{milcVino}
\textbf{Para la próxima sustentación.} Vuelve a esta tabla en seis meses. Verás que las respuestas cambian — no porque lo aprendido cambie, sino porque tú cambias. La evaluación liberadora es una conversación contigo a través del tiempo.
\end{infoband}


\puente{Cerramos con 3 voces sobre por qué la calidad de la pregunta determina la calidad de la respuesta.}

\begin{titlebox}{milcGradeDark}
Triángulo de pensamiento · cierre filosófico
\end{titlebox}

\begin{softbox}{milcVino}{milcGris}{Dussel · lente del nosotros}
\textit{````Quien sabe formular bien una pregunta, tiene la mitad de la respuesta. La palabra precisa es semilla de pensamiento claro.''''}

\textbf{La habilidad de preguntar bien no es nueva}: doña Mercedes la enseñaba en la vereda La Plata. Los abogados la aprenden en su carrera. Los médicos en sus consultas. Los periodistas en sus entrevistas. \emph{Tú la aprendes con la IA, pero te servirá para toda la vida con humanos también}. Quien aprende prompting profundamente aprende también \emph{a pensar con precisión}, lo cual es virtud universal.

\textbf{¿Estoy aprendiendo solo a hablar con IA, o también a pensar con más precisión en todas mis conversaciones?}
\end{softbox}
\linead{\linewidth}\\[1mm]
\linead{\linewidth}

\begin{softbox}{milcMostaza}{milcGris}{Estoicismo · Marco Aurelio · cuidado interior}
\textit{````Pregunta lo que necesitas saber, no lo que crees que debes preguntar. La diferencia es respeto por tu propio tiempo y el del que responde.''''}

\textbf{Mucha gente pregunta lo que ``debería preguntar''} (para parecer culta, para no ofender) en lugar de \emph{lo que realmente necesita saber}. Los LLMs no juzgan: pregunta lo específico. Si necesitas saber qué hacer con tu mamá que se enojó, pregúntalo. Si necesitas ayuda con álgebra de 7°, pídelo en ese nivel. \emph{La autenticidad en el prompt da autenticidad en la respuesta}.

\textbf{¿Pregunto lo que realmente necesito o lo que creo que ``está bien preguntar''?}
\end{softbox}
\linead{\linewidth}\\[1mm]
\linead{\linewidth}

\begin{softbox}{milcTurquesa}{milcGris}{Floridi · lente de la infoesfera}
\textit{````El prompting es la habilidad alfabetizada del siglo XXI. Quien sabe pedir bien, obtiene mucho. Quien no, queda con respuestas pobres.''''}

\textbf{En 2030, saber prompting bien será como hoy saber buscar en Google}: habilidad básica del ciudadano digital. Tu generación llega en momento clave: las apps con IA conversacional se vuelven ubicuas. \emph{Aprender prompting con criterio en grado 7 te coloca en ventaja sobre adultos que recién están empezando}. No es exagerar: tus padres probablemente todavía no dominan esto. Tú sí lo harás.

\textbf{¿Estoy entrenando una habilidad que en 5 años me dará ventaja real, o me limito a usar IA como entretenimiento?}
\end{softbox}
\linead{\linewidth}\\[1mm]
\linead{\linewidth}

\begin{titlebox}{milcVino}
Compromiso final
\end{titlebox}

\small
\begin{tabularx}{\textwidth}{>{\RaggedRight\arraybackslash}p{3.0cm}YYY}
\rowcolor{milcVino}\color{white}\bfseries Criterio & \color{white}\bfseries Lo logré & \color{white}\bfseries En proceso & \color{white}\bfseries Necesito apoyo\\
Comprensión del tema & Explico ideas centrales con mis palabras. & Reconozco ideas pero falta precisión. & Necesito repasar conceptos base.\\
Pensamiento computacional & Apliqué los 4 pilares al diseñar mi producto. & Apliqué algunos pilares. & Necesito releer la fase 2.\\
Aplicación y producto & Mi producto cumple los criterios y se puede revisar. & Producto funcional con ajustes pendientes. & Aún en construcción.\\
Reflexión 5 dimensiones & Respondí cada dimensión con honestidad. & Respondí algunas con detalle. & Mis respuestas son muy generales.\\
Triángulo de pensamiento & Conecté las 3 voces con mi trabajo real. & Conecté al menos una. & No relaciono las voces con mi proceso.\\
\end{tabularx}
\normalsize

\textbf{Hoy entendí que:}\\[1mm]
\linead{\linewidth}\\[1mm]
\linead{\linewidth}

\textbf{Mi compromiso para la próxima sesión es:}\\[1mm]
\linead{\linewidth}\\[1mm]
\linead{\linewidth}

\begin{softbox}{milcMostaza}{milcGris}{Cita de despedida · Marco Aurelio}
\textit{``El obstáculo en el camino se vuelve el camino. Lo que estorba al camino se vuelve camino.''}\\[1mm]
Cada error de hoy, bien observado, se convierte en pieza del camino que pisarás mañana.
\end{softbox}

\small
\begin{tabularx}{\textwidth}{>{\bfseries}p{3.6cm}Y}
\rowcolor{milcGradeDark!12}Nombre completo: & \linead{10cm}\\
Fecha: & \linead{4cm} \quad Curso/grupo: \linead{4cm}\\
Mi firma: & \linead{9cm}\\
\end{tabularx}
\normalsize

\begin{infoband}{milcGradeDark}
\textbf{Para la próxima sesión.} Esta semana, antes de cualquier prompt, piensa CTRF mentalmente. Verás que tus respuestas mejoran de inmediato. La próxima clase aprendes \textbf{prompting avanzado}: técnicas como roles, few-shot, chain-of-thought. Vas a llevar tu prompting al nivel siguiente.
\end{infoband}

\end{document}
